\chapter{Linear Matter Growth and First Structure Observables}

Linear density growth supplies the first late-time structure observable.  The continuum growth equation and its projector discretization are compared in Einstein--de Sitter and flat \(\Lambda\)CDM backgrounds.

\allowdisplaybreaks

\section{Subhorizon dust system}
For pressureless matter,
\begin{align}
\dot\delta&=-\frac{\theta}{a},\\
\dot\theta+H\theta&=-\frac{1}{a}\nabla^2\Phi,\\
\nabla^2\Phi&=4\pi Ga^2\bar\rho_m\delta.
\end{align}
Differentiate the continuity equation:
\begin{align}
\ddot\delta
&=-\frac{\dot\theta}{a}
+\frac{H\theta}{a}.
\end{align}
Using the Euler equation in the form
\begin{equation}
\dot\theta=-H\theta-\frac1a\nabla^2\Phi
\end{equation}
gives
\begin{align}
\ddot\delta
&=\frac{H\theta}{a}
+\frac{1}{a^2}\nabla^2\Phi
+\frac{H\theta}{a}\\
&=-2H\dot\delta
+\frac{1}{a^2}\nabla^2\Phi.
\end{align}
The Poisson equation then yields
\begin{equation}
\frac{1}{a^2}\nabla^2\Phi
=4\pi G\bar\rho_m\delta,
\end{equation}
so elimination of \(\theta\) and \(\Phi\) gives
\begin{equation}
\boxed{
\ddot\delta+2H\dot\delta-4\pi G\bar\rho_m\delta=0.
}
\end{equation}
For the growth factor \(D(a)\),
\begin{align}
\dot D&=aH D_{,a},\\
\ddot D
&=a^2H^2D_{,aa}
+\left(aH^2+a^2HH_{,a}\right)D_{,a}.
\end{align}
Substitution into the time-domain growth equation gives
\begin{equation}
a^2H^2D_{,aa}
+\left(3aH^2+a^2HH_{,a}\right)D_{,a}
-4\pi G\bar\rho_mD=0.
\end{equation}
Using
\begin{equation}
4\pi G\bar\rho_m
=\frac32\Omega_{m0}H_0^2a^{-3}
\end{equation}
and dividing by \(a^2H^2\) yields
\begin{equation}
\boxed{
D_{,aa}
+\left(\frac3a+\frac{H_{,a}}{H}\right)D_{,a}
-\frac{3\Omega_{m0}H_0^2}{2a^5H(a)^2}D=0.
}
\end{equation}
With \(N=\log a\),
\begin{align}
D_{,a}&=\frac1aD_{,N},\\
D_{,aa}&=\frac1{a^2}
\left(D_{,NN}-D_{,N}\right),\\
\frac{aH_{,a}}{H}
&=\frac{\dd\log H}{\dd N},\\
\Omega_m(a)
&=\frac{\Omega_{m0}H_0^2a^{-3}}{H(a)^2}.
\end{align}
Multiplying the \(a\)-equation by \(a^2\) and inserting these identities gives
\begin{equation}
\boxed{
D_{,NN}
+\left[2+\frac{\dd\log H}{\dd N}\right]D_{,N}
-\frac32\Omega_m(a)D=0.
}
\end{equation}

\section{Einstein--de Sitter benchmark}
For
\begin{equation}
\Omega_m(a)=1,
\qquad
\frac{\dd\log H}{\dd N}=-\frac32,
\end{equation}
a power law \(D\propto a^r=\ee^{rN}\) gives the characteristic equation
\begin{align}
r^2+\frac12r-\frac32&=0,\\
(r-1)\left(r+\frac32\right)&=0.
\end{align}
Hence
\begin{equation}
\boxed{
D_+(a)=a,
\qquad
D_-(a)=a^{-3/2}.
}
\end{equation}
The midpoint solver converges with
\begin{equation}
p_{\mathrm{EdS}}=1.968111.
\end{equation}

\section[Flat LambdaCDM growth]{Flat \(\Lambda\)CDM growth}
Let
\begin{align}
E(a)^2&=\Omega_{m0}a^{-3}+1-\Omega_{m0},\\
\Omega_m(a)&=\frac{\Omega_{m0}a^{-3}}{E(a)^2}.
\end{align}
Define
\begin{equation}
I(a):=\int_0^a\frac{\dd a'}{a'^3E(a')^3},
\qquad
I_{,a}=\frac1{a^3E(a)^3}.
\end{equation}
Reduction of order about the decaying solution proportional to \(E(a)\) gives a second solution proportional to \(E(a)I(a)\).  Normalizing it to the matter-era behavior \(D_+\sim a\) fixes the coefficient, so
\begin{equation}
\boxed{
D_+(a)
=\frac{5\Omega_{m0}}{2}
E(a)
\int_0^a\frac{\dd a'}{a'^3E(a')^3}.
}
\end{equation}
For \(\Omega_{m0}=0.3\), the growth ODE and integral solution agree with trajectory residual
\begin{equation}
\varepsilon_{\mathrm{Heath}}
=1.035117\times10^{-12}.
\end{equation}
The midpoint convergence order is
\begin{equation}
p_{\Lambda\mathrm{CDM}}=1.970169.
\end{equation}

\section{Consistent projector Laplacian}
Let
\begin{equation}
\Delta_h\ee^{\ii\bm k\cdot\bm x}
=-\widehat k_h^2\ee^{\ii\bm k\cdot\bm x}.
\end{equation}
The lattice Poisson equation is
\begin{equation}
-\widehat k_h^2\Phi_h
=4\pi Ga^2\bar\rho_m\delta_h.
\end{equation}
The lattice Euler force is
\begin{equation}
\dot\theta_h+H\theta_h
=\frac{\widehat k_h^2}{a}\Phi_h.
\end{equation}
Therefore
\begin{equation}
\boxed{
\frac{\widehat k_h^2}{a}\Phi_h
=-4\pi Ga\bar\rho_m\delta_h.
}
\end{equation}
Differentiating \(\dot\delta_h=-\theta_h/a\) gives
\begin{align}
\ddot\delta_h
&=-\frac{\dot\theta_h}{a}
+\frac{H\theta_h}{a}\\
&=-2H\dot\delta_h
-\frac{\widehat k_h^2}{a^2}\Phi_h.
\end{align}
The discrete Poisson relation now removes the same symbol that appears in the Euler force:
\begin{equation}
-\frac{\widehat k_h^2}{a^2}\Phi_h
=4\pi G\bar\rho_m\delta_h.
\end{equation}
Hence the common projector Laplacian cancels exactly:
\begin{equation}
\boxed{
\ddot\delta_h+2H\dot\delta_h
-4\pi G\bar\rho_m\delta_h=0.
}
\end{equation}
Thus a consistent Euler--Poisson pair introduces no artificial scale dependence in pressureless linear growth.

\section{Growth observables}
Define
\begin{align}
f(a)
&:=\frac{\dd\log D}{\dd\log a}
=\frac{(1/D)\dd D}{\dd N}
=\frac{D_{,N}}{D},\\
D_N(a)&:=\frac{D(a)}{D(1)},\\
\sigma_8(a)&=D_N(a)\sigma_{8,0},\\
f\sigma_8(a)&:=f(a)\sigma_8(a)
=f(a)D_N(a)\sigma_{8,0}.
\end{align}
For
\begin{equation}
\Omega_{m0}=0.3,
\qquad
\sigma_{8,0}=0.8,
\end{equation}
the controlled benchmark gives
\begin{align}
f(1)&=0.512796,\\
f\sigma_8(1)&=0.410237.
\end{align}
These values are conditional on the stated background parameters and normalization, not an observational fit.

\section{Open extensions}
\begin{nextsteps}
\item Add radiation, baryon, and neutrino transfer functions to the projector growth system.
\item Extend the evolution through the superhorizon regime and include free-streaming species.
\item Normalize the matter power spectrum consistently from the primordial sector.
\item Map the predicted growth history to survey observables and a reproducible likelihood.
\item Develop the nonlinear fluid or \(N\)-body limit while preserving the projector constraints.
\end{nextsteps}
